\chapter{الگوریتم‌های کلاسیک}%
\label{ch:classic-algorithms}

جعبه‌ابزارمان کامل است: متغیر، عملگر، شرط، حلقه، تابع و آرایه. در این فصل
با آن‌ها چند الگوریتمِ کلاسیک می‌سازیم؛ الگوریتم‌هایی که دهه‌هاست در قلبِ
درسِ برنامه‌نویسی جا خوش کرده‌اند، چون هم به‌راستی پرکاربردند و هم هر
کدام درسی از «فکرِ الگوریتمی» می‌دهند.

\section{جست‌وجوی خطی}

مسئله: آیا عددِ موردِ نظرِ ما در آرایه هست؟ کجاست؟

سرراست‌ترین راه همان است که خودمان در زندگی می‌کنیم: از اول یکی‌یکی نگاه
کن تا یا پیدایش کنی یا فهرست تمام شود. نامِ این روش \textbf{جست‌وجوی خطی}
است. قرارداد این است که اگر پیدا شد، \emph{اندیسش} را برگردانیم و اگر نه،
عددِ \code{-1} را (که چون اندیسِ معتبری نیست، پیام‌آورِ خوبی برای «نبود»
است):

\begin{salam}
تابع جستجوی خطی(اعداد : صحیح[:], هدف : صحیح) : صحیح:
    متغیر i : صحیح = 0
    تاوقتی i < len(اعداد):
        اگر اعداد[i] == هدف:
            بازگشت i
        تمام
        i = i + 1
    تمام
    بازگشت -1
تمام

تابع اصلی:
    کدها : صحیح[6] = [40, 12, 73, 25, 66, 12]
    چاپ جستجوی خطی(کدها, 25)
    چاپ جستجوی خطی(کدها, 99)
    چاپ جستجوی خطی(کدها, 12)
تمام
\end{salam}

\begin{progout}
3
-1
1
\end{progout}

سه فراخوانی، سه درس: هدفِ موجود اندیسش را می‌دهد؛ هدفِ ناموجود \code{-1}؛
و هدفِ تکراری اندیسِ \emph{اولین} حضورش را، چون \code{بازگشت} در همان
لحظهٔ پیدا شدن تابع را ترک می‌کند و به ادامهٔ گشتن نمی‌رسد.

\begin{deep}[نگاهی به درون: هزینهٔ الگوریتم]
برای آرایه‌ای با هزار خانه، جست‌وجوی خطی در بدترین حالت (وقتی هدف نیست)
هزار مقایسه می‌کند؛ برای یک میلیون خانه، یک میلیون. می‌گوییم هزینه‌اش
«متناسب با طولِ ورودی» است. سنجیدنِ الگوریتم‌ها با این عینک (اگر ورودی ده
برابر شود، کار چند برابر می‌شود؟) شاخهٔ مهمی از دانشِ رایانه است. برای
کنجکاوها: اگر آرایه از قبل مرتب باشد، می‌توان با روشِ «جست‌وجوی دودویی»
(هر بار وسط را نگاه کن و نیمهٔ بی‌ربط را دور بریز) میلیون خانه را با
حدودِ بیست مقایسه گشت! پیاده‌سازی‌اش تمرینِ پایانیِ همین فصل است.
\end{deep}

\section{مرتب‌سازیِ حبابی}

مسئله: آرایه را از کوچک به بزرگ بچین.

مرتب‌سازی شاید پرمطالعه‌ترین مسئلهٔ تاریخِ رایانه باشد و ده‌ها الگوریتم
دارد. ساده‌ترینشان برای یادگیری، \textbf{مرتب‌سازیِ حبابی} است با یک
ایدهٔ کوچک: از اولِ آرایه راه بیفت و هر دو همسایه که ترتیبشان غلط است را
جابه‌جا کن؛ این گشت را آن‌قدر تکرار کن تا آرایه مرتب شود. در هر گشت،
بزرگ‌ترین عددِ باقی‌مانده مثلِ حباب به انتهای آرایه می‌رسد (نامش از
همین‌جاست).

اول ابزارِ کوچکش: جابه‌جاییِ دو خانه. یادتان هست در تمرینِ فصلِ
\ref{ch:variables} برای جابه‌جا کردنِ دو متغیر به یک جای امنِ سوم نیاز
داشتیم؟ همان نقشه اینجا کارِ اصلی را می‌کند:

\begin{salam}
تابع مرتب سازی حبابی(اعداد : صحیح[:]):
    ن : صحیح = len(اعداد)
    تکرار ن - 1:
        متغیر i : صحیح = 0
        تاوقتی i < ن - 1:
            اگر اعداد[i] > اعداد[i + 1]:
                موقت : صحیح = اعداد[i]
                اعداد[i] = اعداد[i + 1]
                اعداد[i + 1] = موقت
            تمام
            i = i + 1
        تمام
    تمام
تمام

تابع اصلی:
    متغیر نمرات : صحیح[6] = [15, 3, 18, 7, 12, 9]
    مرتب سازی حبابی(نمرات)
    هر نمره در نمرات:
        بنویس نمره
        بنویس " "
    تمام
    چاپ ""
تمام
\end{salam}

\begin{progout}
3 7 9 12 15 18
\end{progout}

ساختار، همان حلقهٔ تودرتوی فصلِ \ref{ch:loops} است: حلقهٔ درونی یک «گشتِ»
کامل روی همسایه‌هاست و حلقهٔ بیرونی این گشت را \code{ن - 1} بار تکرار
می‌کند، چون هر گشت دستِ‌کم یک عدد را سرِ جای نهایی‌اش می‌نشاند. برای
فهمِ عمیق، یک اجرای دستیِ کوچک روی آرایهٔ \code{[3, 1, 2]} انجام دهید:
گشتِ اول \code{[1, 2, 3]} می‌سازد (دو جابه‌جایی) و گشتِ دوم دیگر چیزی
جابه‌جا نمی‌کند. مرتب شد!

\begin{note}
چرا مرتب‌سازی این‌قدر مهم است؟ چون دادهٔ مرتب، درِ الگوریتم‌های سریع‌تر
را باز می‌کند (مثلِ همان جست‌وجوی دودویی) و خروجیِ مرتب برای انسان
خواناتر است. دفترچه‌تلفن را تصور کنید اگر نام‌ها به همان ترتیبِ ثبت‌نام
چاپ شده بودند!
\end{note}

\section{بزرگ‌ترین مقسوم‌علیهِ مشترک: الگوریتمِ اقلیدس}

مسئله: بزرگ‌ترین عددی که دو عددِ داده‌شده هر دو بر آن بخش‌پذیرند
(ب.م.م) را بیاب؛ همان که در ساده کردنِ کسرها به کار می‌رود.

راهِ ساده‌لوحانه این است که از عددِ کوچک‌تر رو به پایین همهٔ اعداد را
امتحان کنیم. کار می‌کند اما کُند است. حدودِ ۲۳ قرن پیش اقلیدس راهی
درخشان یافت که هنوز هم بهترین است و بر این حقیقت تکیه دارد: ب.م.مِ دو
عدد، با ب.م.مِ «عددِ کوچک‌تر» و «باقی‌ماندهٔ تقسیمِ بزرگ‌تر بر کوچک‌تر»
برابر است؛ و هر گاه یکی از دو عدد صفر شود، آن دیگری جواب است.

\begin{salam}
تابع ب م م(الف : صحیح, ب : صحیح) : صحیح:
    متغیر اول : صحیح = الف
    متغیر دوم : صحیح = ب
    تاوقتی دوم != 0:
        باقیمانده : صحیح = اول % دوم
        اول = دوم
        دوم = باقیمانده
    تمام
    بازگشت اول
تمام

تابع اصلی:
    چاپ ب م م(48, 36)
    چاپ ب م م(270, 192)
    چاپ ب م م(17, 5)
تمام
\end{salam}

\begin{progout}
12
6
1
\end{progout}

اجرای دستی برای \code{(48, 36)}: باقی‌ماندهٔ \code{48 \% 36} می‌شود ۱۲،
پس جفتِ تازه \code{(36, 12)} است؛ باقی‌ماندهٔ \code{36 \% 12} صفر است،
پس جفت \code{(12, 0)} می‌شود و حلقه می‌ایستد: جواب ۱۲. فقط دو دور! زیباییِ
الگوریتمِ خوب همین است: نه کدِ بیشتر، \emph{فکرِ} بهتر.

\section{دنبالهٔ فیبوناچی}

دنباله‌ای مشهور که با ۰ و ۱ شروع می‌شود و هر جملهٔ بعدی مجموعِ دو جملهٔ
قبلی است: ۰، ۱، ۱، ۲، ۳، ۵، ۸، ۱۳ و همین‌طور تا همیشه. این دنباله در
طبیعت (آرایشِ دانه‌های آفتاب‌گردان، مارپیچِ صدف‌ها) سر و کله‌اش پیدا
می‌شود و در دنیای الگوریتم هم مثالِ محبوبِ «ساختنِ هر جمله از جمله‌های
پیشین» است. برای تولیدش فقط باید دو جملهٔ آخر را به خاطر داشته باشیم؛
دو متغیر و همان رقصِ جابه‌جایی:

\begin{salam}
تابع اصلی:
    متغیر قبلی : صحیح = 0
    متغیر فعلی : صحیح = 1
    تکرار 20:
        بنویس قبلی
        بنویس " "
        بعدی : صحیح = قبلی + فعلی
        قبلی = فعلی
        فعلی = بعدی
    تمام
    چاپ ""
تمام
\end{salam}

\begin{progout}
0 1 1 2 3 5 8 13 21 34 55 89 144 233 377 610 987 1597 2584 4181
\end{progout}

\section{وارونه کردنِ یک عدد}

مسئلهٔ پایانی: عددی بگیر و وارونه‌اش را بساز؛ \code{1404} بشود
\code{4041}. ایدهٔ حل، ترکیبِ دو ترفندِ آشنای فصلِ \ref{ch:operators}
است: \code{\% 10} رقمِ یکان را می‌کند و \code{/ 10} آن را می‌اندازد. هر
رقمی که کندیم را از راست به عددِ وارونه می‌چسبانیم، با ضرب در ۱۰:

\begin{salam}
تابع وارونه(عدد : صحیح) : صحیح:
    متغیر مانده : صحیح = عدد
    متغیر نتیجه : صحیح = 0
    تاوقتی مانده > 0:
        رقم : صحیح = مانده % 10
        نتیجه = نتیجه * 10 + رقم
        مانده = مانده / 10
    تمام
    بازگشت نتیجه
تمام

تابع اصلی:
    چاپ وارونه(1404)
    چاپ وارونه(12321)
تمام
\end{salam}

\begin{progout}
4041
12321
\end{progout}

فراخوانیِ دوم چیزِ قشنگی نشان می‌دهد: عددی که با وارونه‌اش برابر است
(\emph{متقارن} یا آینه‌ای). حالا آزمودنِ متقارن بودن یک خط بیشتر نیست:
\code{عدد == وارونه(عدد)}.

\section{تمرین‌ها}

\exercise جست‌وجوی خطی را طوری تغییر دهید که به‌جای اولین حضور، اندیسِ
\emph{آخرین} حضورِ هدف را برگرداند (دو راه دارد: پیمایش از آخر، یا ادامه
دادنِ گشت و به خاطر سپردنِ آخرین جای دیده‌شده؛ هر دو را امتحان کنید).

\exercise مرتب‌سازیِ حبابی را با یک اجرای دستی روی \code{[5, 1, 4, 2]}
دنبال کنید و بعد از هر گشتِ کامل، وضعِ آرایه را بنویسید.

\exercise مرتب‌سازیِ حبابی را از بزرگ به کوچک بازنویسی کنید. فقط یک نویسه
باید عوض شود؛ کدام؟

\exercise با کمکِ تابعِ \code{ب م م}، تابعِ
\code{ک م م(الف صحیح, ب صحیح) صحیح} (کوچک‌ترین مضربِ مشترک) را بنویسید.
راهنمایی: حاصل‌ضربِ دو عدد برابر است با حاصل‌ضربِ ب.م.م و ک.م.مِ آن‌ها.

\exercise تابعی بنویسید که مجموعِ رقم‌های یک عدد را برگرداند
(\code{1404} بدهد \code{9}). اسکلتش همان تابعِ \code{وارونه} است.

\exercise (چالشی) جست‌وجوی دودویی: تابعی بنویسید که در آرایه‌ای
\emph{مرتب} بگردد. دو نشانگرِ \code{پایین} و \code{بالا} نگه دارید که اول
به دو سرِ آرایه اشاره می‌کنند؛ در هر دور خانهٔ وسط را ببینید: اگر هدف
بود تمام، اگر هدف کوچک‌تر بود \code{بالا} را به وسط منتقل کنید و اگر
بزرگ‌تر بود \code{پایین} را؛ تا وقتی \code{پایین <= بالا}. با آرایهٔ
مرتبِ خروجیِ مرتب‌سازیِ حبابی بیازماییدش.
